%==============================================================================
\appsection{Full Covariance Derivation}
\label{app:full_cov}
%==============================================================================

In the main text, I derived hyperpriors for a model with $p$ independent
parameters, each governed by its own Normal-Inverse-Gamma hyperprior. Here I
extend the derivation to the case where the $p$ parameters have a full
covariance structure, using Normal-Inverse-Wishart (NIW) hyperpriors. This
extension is relevant when the researcher believes that the population-level
parameters may be correlated---for example, that subjects with higher risk
aversion also tend to have lower probability weighting curvature.

\subsection{Setup}

Suppose the direct prior for a single individual is multivariate normal:
\begin{equation}
  \label{eq:mv_direct_prior}
  \bth_i \sim \text{MVN}(\boldsymbol{\mu}^*, \boldsymbol{\Sigma}^*)
\end{equation}
where $\bth_i \in \R^p$ is the parameter vector,
$\boldsymbol{\mu}^* \in \R^p$ is the direct prior mean, and
$\boldsymbol{\Sigma}^*$ is a $p \times p$ positive definite covariance matrix.
Note that when $\boldsymbol{\Sigma}^*$ is diagonal, this reduces to $p$
independent normal distributions, recovering the setting of
Section~\ref{sec:variance_decomposition}.

\subsection{Normal-Inverse-Wishart Hierarchical Structure}

For the hierarchical model with correlated parameters, a natural conjugate
choice is the Normal-Inverse-Wishart (NIW) hyperprior:
\begin{align}
  \label{eq:niw_mu}
  \boldsymbol{\mu} \mid \boldsymbol{\Sigma} &\sim
    \text{MVN}\!\left(\boldsymbol{\mu}_0,\;
    \frac{\boldsymbol{\Sigma}}{\kappa_0}\right) \\[4pt]
  \label{eq:niw_sigma}
  \boldsymbol{\Sigma} &\sim \text{Inverse-Wishart}(\boldsymbol{\Psi}_0, \nu_0)
\end{align}
where $\boldsymbol{\mu}_0 \in \R^p$ is the prior mean for the
population mean, $\kappa_0 > 0$ scales the variance of the population mean
relative to the population covariance, $\boldsymbol{\Psi}_0$ is the $p \times
p$ positive definite scale matrix, and $\nu_0 > p + 1$ is the degrees of
freedom parameter.\footnote{
  The NIW scaling parameter $\kappa_0$ is standard notation in the conjugate
  prior literature and should not be confused with the coefficient of variation
  of heterogeneity $\kappa$ introduced in
  Section~\ref{sec:variance_decomposition}.
}

\subsection{Mean Constraint}

By the law of iterated expectations:
\begin{equation}
  \label{eq:niw_mean_constraint}
  \boldsymbol{\mu}_0 = \boldsymbol{\mu}^*
\end{equation}

\subsection{Covariance Constraint}

The matrix analog of the law of total variance gives:
\begin{equation}
  \label{eq:matrix_total_var}
  \text{Cov}[\bth_i]
    = \E{\text{Cov}[\bth_i \mid \boldsymbol{\mu}, \boldsymbol{\Sigma}]}
    + \text{Cov}[\E{\bth_i \mid \boldsymbol{\mu}}]
\end{equation}

The unconditional covariance is:
\begin{equation}
  \label{eq:niw_total_cov}
  \boldsymbol{\Sigma}^*
    = \E{\boldsymbol{\Sigma}} + \text{Cov}[\boldsymbol{\mu}]
    = \frac{\boldsymbol{\Psi}_0}{\nu_0 - p - 1}
      \cdot \frac{\kappa_0 + 1}{\kappa_0}
\end{equation}

\subsection{Variance Partition with \texorpdfstring{$\rho$}{rho}}

The $\rho$ decomposition applies identically:
\begin{align}
  \label{eq:niw_cov_mu}
  \text{Cov}[\boldsymbol{\mu}]
    &= \rho \cdot \boldsymbol{\Sigma}^* \\[4pt]
  \label{eq:niw_E_sigma}
  \E{\boldsymbol{\Sigma}}
    &= (1 - \rho) \cdot \boldsymbol{\Sigma}^*
\end{align}

\subsection{Derived Hyperparameters}

The researcher specifies a direct prior $\boldsymbol{\mu}^*$,
$\boldsymbol{\Sigma}^*$, the variance partition ratio $\rho \in (0, 1)$, and
the degrees of freedom $\nu_0 > p + 1$. The NIW hyperprior parameters are then:
\begin{align}
  \boldsymbol{\mu}_0 &= \boldsymbol{\mu}^*
    \label{eq:niw_hp_mu} \\[3pt]
  \kappa_0 &= \frac{1 - \rho}{\rho}
    \label{eq:niw_hp_kappa} \\[3pt]
  \boldsymbol{\Psi}_0 &= (1-\rho) \cdot \boldsymbol{\Sigma}^* \cdot (\nu_0 - p - 1)
    \label{eq:niw_hp_psi} \\[3pt]
  \nu_0 &= \nu_0 \quad \text{(researcher specified)}
    \label{eq:niw_hp_nu}
\end{align}

\subsection{Verification}

Substituting the derived hyperparameters back into (\ref{eq:niw_total_cov}):
\begin{align}
  \E{\boldsymbol{\Sigma}} + \text{Cov}[\boldsymbol{\mu}]
  &= \frac{(1-\rho) \boldsymbol{\Sigma}^* (\nu_0 - p - 1)}{\nu_0 - p - 1}
     \left(1 + \frac{\rho}{1 - \rho}\right) = \boldsymbol{\Sigma}^*
\end{align}

\subsection{Relationship to the Independent Case}

When $\boldsymbol{\Sigma}^* = \text{diag}(\sigma_1^{2*}, \ldots,
\sigma_p^{2*})$, the Inverse-Wishart reduces to $p$ independent Inverse-Gamma
distributions:
\begin{equation}
  \Sigma_{jj} \sim \text{Inverse-Gamma}\!\left(
    \frac{\nu_0 - p + 1}{2},\;
    \frac{(1-\rho) \sigma_j^{2*} (\nu_0 - p - 1)}{2}
  \right)
\end{equation}
recovering the independent-parameter formulas with $\alpha_0 = (\nu_0 - p + 1)/2$.

\subsection{Computational Note}

For MCMC sampling, the covariance matrix $\boldsymbol{\Sigma}$ is
parameterized via its Cholesky decomposition $\boldsymbol{\Sigma} =
\mathbf{L}\mathbf{L}^\top$, where $\mathbf{L}$ is lower triangular with
positive diagonal stored on the log scale.

\pagebreak
